\subsection{Affine structure and linear thresholds}\label{sec:affine}

Parity checks, linear hashes and CRCs make rare events with $p=2^{-\mathrm{rank}}$,
which Gaussian elimination computes exactly under the uniform distribution. They are
also where decomposition fails. Take $m\le n/6$ random parity checks and an event
inside their solution set that holds on at least a $1/\poly(n)$ fraction of it. With
high probability over the checks, the envelope estimator with decomposable
envelopes whose leaves have at most $2n/3$ variables, context-width compilation and the
node budget behind DFSB's completeness
guarantee then cost Monte Carlo's price up to a polynomial factor, and every use of a
DNNF envelope costs at least its square root
(\cref{thm:parity-barrier,cor:barrier-methods}). This subsection
makes affine structure usable soundly, proves that barrier, and maps where
tractability ends: under biased product distributions and in combination with
thresholds. We write $U=\{i:q_i=1/2\}$ for the uniform coordinates of $D$ and
$D_i(b)=\Pr_D[x_i=b]$.

\paragraph{Certified affine gates.}
The algebraic normal form (ANF) of a Boolean function is its unique representation
as a multilinear polynomial over $\F_2$. We propagate ANFs through the circuit
($\wedge\mapsto ab$ with $x_i^2=x_i$, $\vee\mapsto a+b+ab$, $\neg\mapsto1+a$,
$\oplus\mapsto a+b$), store the linear part as a bit vector, and abort a gate when
its nonlinear part exceeds $T$ monomials; threshold gates are opaque. A gate whose
ANF is computed is computed exactly, because the rules are identities and the ANF is
unique; if its degree is at most one the gate is \emph{certified affine},
$g(x)=a_g\cdot x+c_g$ for all $x$. A chain of two-input XOR gadgets
$(a\vee b)\wedge\neg(a\wedge b)$ on $k$ inputs is certified with $T=k-1$ and $O(k^2)$
work per gate (our implementation also caps the work of each product); a balanced
tree of such gadgets needs $T\ge\lfloor k/2\rfloor\lceil k/2\rceil$, because its last
$\vee$ multiplies two disjoint linear forms. A gate is in \emph{solved form} on a coordinate $i$,
$g(x)=x_i\oplus G(x_{-i})$, if the ANF computed with $x_i$ as one symbol and every
gate whose support avoids $i$ as an opaque symbol contains $x_i$ only in the
monomial $x_i$. Substituting Boolean functions for symbols is a homomorphism of
Boolean rings, so this is a certificate as well, with $G=g|_{x_i=0}$. Certification
has to be deductive. A gate computing $\ell\oplus\one_{x^*}$ agrees with the affine
form $\ell$ except at a hidden point, which $T$ evaluations miss with probability
$1-T2^{-n}$; yet in $C=[g=1]\wedge[Ax=b]$ with $\ell\equiv0$ on $\{Ax=b\}\ni x^*$ it
decides between $p=2^{-n}$ and $p=0$ (compare \cref{thm:blackbox}).

\paragraph{Exact presolve.}
The \emph{$\wedge$-spine} of $C$ is the set of pairs $(g,b)$ reached from
$(\mathrm{out},1)$ through $\wedge$-gates with $b=1$, $\vee$-gates with $b=0$ and
negations, without entering certified affine or solved-form gates; then
$C(x)=\bigwedge_{(g,b)}[g(x)=b]$ identically.

\begin{proposition}[exact affine and triangular presolve]\label{prop:presolve}
Suppose $C(x)=\bigwedge_{j\le k}[x_{p_j}=\pi_j(x)]\wedge R(x)$ for all $x$, where the
$p_j$ are distinct and $\pi_j$ does not read $x_{p_j},\dots,x_{p_k}$. Let
$P=\{p_1,\dots,p_k\}$, and for $x\in\bits^n$ let $\tilde x$ agree with $x$ off $P$
and satisfy $\tilde x_{p_j}=\pi_j(\tilde x)$, fixed for $j=1,\dots,k$ in turn; let
$\tilde R(x)=R(\tilde x)$, which does not read $P$. Then, for every product
distribution $D$,
\[
p=M\,\E_D[W],\qquad
W(x)=\tilde R(x)\prod_{j\le k}\frac{D_{p_j}(\tilde x_{p_j})}{\max(q_{p_j},1-q_{p_j})}\in[0,1],\qquad
M=\prod_{j\le k}\max(q_{p_j},1-q_{p_j}).
\]
In particular $p=2^{-k}D(\tilde R)$ if $P\subseteq U$; and if
$|q_{p_j}-\frac12|\le\eta$ for all $j$, then $\E_D W\ge\rho^{-k}D(\tilde R)$ with
$\rho=\frac{1+2\eta}{1-2\eta}$.
\end{proposition}

\begin{proof}
For fixed $x_{-P}$ the constraints determine $x_{p_1},\dots,x_{p_k}$ in turn, since
$\pi_j$ reads only coordinates off $P$ and earlier pivots. So exactly one completion,
$\tilde x$, satisfies them, and
$p=\sum_{x_{-P}}D_{-P}(x_{-P})D_P(\tilde x_P)R(\tilde x)=M\,\E_DW$, as $W$ does not
depend on $x_P$. If $P\subseteq U$, then $D_P\equiv2^{-k}=M$ and $W=\tilde R$; with
$|q_{p_j}-\frac12|\le\eta$ every ratio in $W$ is at least $\rho^{-1}$.
\end{proof}

Two kinds of spine conjuncts have this form. Certified affine conjuncts
$[a_g\cdot x=b\oplus c_g]$ of rank $r$, in reduced row echelon form
$x_P=b'+A'x_{[n]\setminus P}$, have every $\pi_j$ reading only non-pivots; we choose
pivots in $U$ and, when possible, outside the support of the remaining conjuncts. A
solved-form conjunct $[x_i\oplus G(x_{-i})=b]$ gives $\pi=b\oplus G$; the presolve
adds such conjuncts greedily, keeping an order in which no earlier $G$ and no affine
row reads a later pivot. It then builds $\tilde R$ as a circuit (the pivots are
replaced by $\oplus$-gates of free inputs and by copies of the cones of the $G$'s) and
returns $M$ times any sound interval for $\E_DW$: by default the KL procedure of
\cref{thm:envelope}(iv) on $W$, and, when $P\subseteq U$, any sound method applied to
$\tilde R$, such as DFSB or the meta-procedure. The certified bound $p\le M$ always
holds. So the presolve is exact and sound on every circuit, explains $k$ bits of
surprise at polynomial cost, and leaves $\log_2(1/D(\tilde R))$ bits to the methods
that follow; under biased pivots each pivot costs a factor up to $\rho$ in samples.
Rule (F) of DFSB also applies Gaussian elimination to affine conjuncts at the root,
but it keeps only the literals the system forces, that is, rows of the reduced system
with a single variable. A rank-$r$ system without such rows explains nothing there,
and $r$ bits here. Examples, under the uniform distribution (or whenever the pivots can
be chosen uniform): the event $[\mathrm{CRC}(x)=0]\wedge R(x)$ has
$p=2^{-r}D(\tilde R)$ with $r$ the rank of the CRC map, and a Feistel round $[L\oplus f(R)=t]$,
whose pivots are the coordinates of $L$, has probability $2^{-|L|}$ for every round
function $f$ (E4 in \cref{tab:affine}). More
generally, under the uniform distribution a map $F$ composed of coordinate
permutations and triangular flips
$\tau(x)=x\oplus g(x_{-i})e_i$ with $i\in U$ preserves $D$ (each flip is a
measure-preserving involution of every fibre $\{x_{-i}=\text{const}\}$), so
$D(R\circ F)=D(R)$ for every $R$, whatever the round functions; our implementation
detects only the one-round case. Gaussian elimination is standard in SAT and
model-counting tools, and eliminating variables defined by others is a known
preprocessing step for model counting \cite{lagniez2016}. What the presolve adds is
its use as an exact surprise-removal step that composes with every sound estimator,
together with the biased-pivot version.

\paragraph{Affine-union envelopes.}
Decomposition fails on parity checks, but abstract interpretation over unions of
affine subspaces succeeds, and the result is an envelope.

\begin{theorem}[affine envelopes]\label{thm:affenv}
For every gate $g$ compute lists $E^+(g),E^-(g)$ of affine subspaces of $\F_2^n$. A
certified affine gate gets $E^+=\{a_g\cdot x=1+c_g\}$ and $E^-=\{a_g\cdot x=c_g\}$;
$\neg$ swaps the lists; for $g=g_1\wedge g_2$ let $E^+(g)$ be the nonempty
intersections $V_1\cap V_2$, $V_j\in E^+(g_j)$, if there are at most $L$ of them, and
otherwise the operand list of smaller mass; let $E^-(g)=E^-(g_1)\cup E^-(g_2)$ if it
has at most $L$ entries, and otherwise $\{\F_2^n\}$; $\vee$ is dual, $\oplus$ uses
$E^+(a\oplus b)=E^+(a)\sqcap E^-(b)\,\cup\,E^-(a)\sqcap E^+(b)$ (and similarly for
$E^-$), wider gates are folded, and other gates get $\{\F_2^n\}$ for both lists.
Duplicate and contained subspaces may be removed.
\begin{enumerate}[label=(\alph*)]
\item For every circuit and every gate, $g^{-1}(1)\subseteq\bigcup E^+(g)$ and
$g^{-1}(0)\subseteq\bigcup E^-(g)$. Hence the function
$\varphi=\sum_{V\in E^+(\mathrm{out})}\one_V$ dominates $C$, and
$E^+(\mathrm{out})=\emptyset$ certifies $p=0$. The lists cost
$\poly(|C|,n,L,T)$ time.
\item If every constraint of $V$ is on uniform coordinates, $D(V)=2^{-\mathrm{codim}V}$
and $D\mid V$ is sampled exactly by drawing $x\sim D$ and re-solving the pivots; for
other $V$, \cref{prop:bw-dp} below gives $D(V)$ and exact samples along a given linear
order of the coordinates, in time exponential in its width (our implementation uses
the input order here, and orders supplied by a prover in \cref{sec:xpt}); if the width
exceeds a cap, $V$ is replaced by its implied constraints on $U$ (a superset). With
these, the envelope estimator (\cref{thm:envelope}) is sound on
every circuit, with $K=\sum_{V\in E^+(\mathrm{out})}D(V)/p$; it is the
Karp--Luby--Madras estimator for the union.
\item If $C=\bigvee_{j\le L}([A_jx=b_j]\wedge R_j)$ with certified affine conjuncts and
opaque $R_j$ (for instance threshold gates), then $K\le L/\min_j\Pr_D[R_j\mid A_jx=b_j]$
under the uniform distribution.
\end{enumerate}
\end{theorem}

\begin{proof}
(a) Induction over the gates. A certified gate equals its affine form. $\neg$ swaps
the one-set and the zero-set. If $g=g_1\wedge g_2$ and $g(x)=1$, then $x$ lies in some
$V_1\cap V_2$, and both coarsenings are supersets; if $g(x)=0$ then $x\in g_1^{-1}(0)\cup
g_2^{-1}(0)$. The other gates are dual or are unions of such cases, and multiplicities
only increase $\varphi$. (b) With constraints on $U$ only, the free uniform coordinates,
the pivots solved from them, and the other coordinates drawn from $D$ give each $x\in
V$ probability $D(x)2^{\mathrm{codim}V}$. The rest is \cref{thm:envelope} with the
integer-valued $\varphi$. (c) Opaque $R_j$ have $E^+=\{\F_2^n\}$, so
$E^+(\mathrm{out})\subseteq\{V_j\}$ with $V_j=\{A_jx=b_j\}$, and
$\sum_jD(V_j)\le L\max_jD(V_j)$ while $p\ge D(V_{j^*})\Pr[R_{j^*}\mid V_{j^*}]$ for the
maximiser $j^*$.
\end{proof}

On the $120$-input instance E1 of \cref{tab:affine} (45 gadget parity checks and a
3-DNF) the envelope has six pieces and $K=1.39$.
Decomposability can be generalised to linear independence, which keeps exact masses
and samplers under the uniform distribution.

\begin{proposition}[XOR-DNNF envelopes]\label{prop:xordnnf}
Let $\Phi$ be an arithmetic circuit with nonnegative weights whose leaves are
nonnegative constants and functions $f(\ell(x))$ of a single nonconstant affine form,
and in which, at every $\Prod$ gate, the spaces $\mathcal L(c)$ spanned by the linear
parts of the leaf forms below the children $c$ form a direct sum. Under the uniform
distribution the bottom-up pass ($\Sum\mapsto$ weighted sum, $\Prod\mapsto$ product,
leaf $\mapsto(f(0)+f(1))/2$) computes $\E\varphi$ exactly, and the following samples
exactly from $D\varphi/\E\varphi$: choose a derivation top-down as in
\cref{lem:envmass}, give each of its leaves the value $v$ with probability
proportional to $f(v)$, and draw $x$ uniformly from the affine subspace on which every
leaf form takes its value.
\end{proposition}

The proof is in \cref{app:affine}. DNNF envelopes (forms $x_i$) and unions of affine
subspaces (products of leaves $[\ell=b]$ with independent forms) are special cases;
$\one_{\{Hx=s\}}$ is an XOR-DNNF with $m$ leaves and $K=1$, while DNNFs whose leaves
have at most $2n/3$ variables need at least $2^{\lambda_{\rm bal}}/(n+2)-2$ wires by the
next theorem (part (b) with $C=E=\one_{\{Hx=s\}}$). Affine decision trees \cite{koriche2013}
use affine decision nodes for exact model counting; XOR-DNNF envelopes also allow
nondeterministic sums and arbitrary leaf functions, at the price of looseness.

\paragraph{The parity barrier for decomposable envelopes.}
For an affine subspace $V\subseteq\F_2^n$ and a cut $(Y,Z)$ of $[n]$ let
$\lambda_V(Y)=\dim\pi_Y(V)+\dim\pi_Z(V)-\dim V$, with $\pi_Y$ the coordinate
projection. For $V=\{Hx=s\}\ne\emptyset$ this is
$\mathrm{rank}\,H_Y+\mathrm{rank}\,H_Z-\mathrm{rank}\,H$, the connectivity function of
the binary matroid of $H$. Let
$\lambda_{\rm bal}(V)=\min\{\lambda_V(Y):n/3<|Y|\le2n/3\}$. Call an envelope circuit or
DNNF \emph{smooth} if every $\Sum$ (or $\vee$) child has the support of its parent
and the root has support $[n]$; leaves may be arbitrary nonnegative (Boolean)
functions of their declared supports. For a circuit $\Phi$ let $G(\Phi)$ be the set
of its nodes $v$ with $n/3<|\sv(v)|\le2n/3$.

\begin{theorem}[parity barrier]\label{thm:parity-barrier}
Let $D$ be uniform, $V$ a nonempty affine subspace, and $C$ a circuit with
$\emptyset\ne C^{-1}(1)\subseteq V$; let $c=|C^{-1}(1)|/|V|=\Pr[C\mid V]$.
\begin{enumerate}[label=(\alph*)]
\item Let $\Phi$ be a smooth envelope circuit with fan-in-two $\Prod$ gates whose
leaves have supports of size at most $2n/3$. If $\varphi\ge1$ on $C^{-1}(1)$ and
$K=\E\varphi/p$, then $|G(\Phi)|\cdot K\ge c\,2^{\lambda_{\rm bal}(V)}$.
\item Let $E$ be a smooth DNNF with fan-in-two $\wedge$-gates whose leaves have
supports of size at most $2n/3$, with $C\Rightarrow E$. Then
$|G(E)|^2\,D(E)/p\ge c\,2^{\lambda_{\rm bal}(V)}$, and $|G(E)|\ge c\,2^{\lambda_{\rm
bal}(V)}$ if moreover $E\Rightarrow[x\in V]$ (for instance if $E\equiv C$).
\end{enumerate}
Every envelope circuit or DNNF with $S$ wires computes the same function as a smooth
one with fan-in-two products and at most $(n+2)(S+2)$ nodes.
\end{theorem}

The proof (\cref{app:affine}) walks down each derivation, following the larger
child at products, to its first node with at most $2n/3$ variables. This writes
$\varphi$ as a sum, over balanced nodes $g$, of products $f_g(x_Y)h_g(x_Z)$. A
Cauchy--Schwarz argument over the $2^{\lambda_V(Y)}$ blocks $A_i\times B_i$ into which
$V$ splits across the cut then shows that such a product covers at most a
$2^{-\lambda_V(Y)/2}$ fraction of $V$ per unit of mass. Part (b) with $E\equiv C$
is the rectangle-cover bound behind known DNNF lower bounds for linear codes
\cite{duris2004,bova2016}. Here the bound is for envelopes, trades size against
looseness, applies to every subset of $V$, and is parametrised by the connectivity
function.

\begin{corollary}\label{cor:parity-ghw}
If $H\in\F_2^{m\times n}$ has rank $m$ and row space $R$, then for every nonempty
$V=\{Hx=s\}$, $\lambda_V(Y)=m-\dim(R\cap\F_2^Y)-\dim(R\cap\F_2^Z)$, where $\F_2^T$
denotes the vectors supported in $T$; hence $\lambda_{\rm bal}(V)\ge m-2j^*$ with
$j^*=\max\{r:d_r(R)\le2n/3\}$ and $d_r$ the $r$-th generalised Hamming weight
\cite{wei1991}.
\end{corollary}

\begin{corollary}[random parity checks]\label{cor:random-parity}
Let $H$ be uniform in $\F_2^{m\times n}$ with $m\le n/6$. With probability at least
$1-2^{1-n/6}$, $H$ has rank $m$ and $\dim(R\cap\F_2^T)\le6$ for every $|T|\le2n/3$;
then $\lambda_{\rm bal}(V)\ge m-12$ for every nonempty $V=\{Hx=s\}$. On this event,
for every circuit $C$ with $C^{-1}(1)\subseteq V$ and $c=\Pr[C\mid V]$ (so that
$p=c2^{-m}$), every $\Phi$ as in (a) has $|G(\Phi)|K\ge c\,2^{m-12}=c^2/(4096p)$ and
every $E$ as in (b) has $|G(E)|^2D(E)/p\ge c^2/(4096p)$.
\end{corollary}

\begin{corollary}[what the barrier rules out]\label{cor:barrier-methods}
On the event of \cref{cor:random-parity}, for $C^{-1}(1)\subseteq V=\{Hx=s\}$:
\begin{enumerate}[label=(\roman*)]
\item \emph{Envelopes.} The envelope estimator with any envelope circuit $\Phi$
whose leaves have supports of at most $2n/3$ variables (for instance literal and
constant leaves, the (R8) atoms of \cref{sec:rules} when $a\le2n/3$, (R7) boxes on at
most $2n/3$ variables, and the leaves of \cref{thm:cw}) evaluates $\varphi$ at every sample and needs $\Omega(K)$ samples, so its work is
$\Omega(c^2/(np))$: Monte Carlo's number of samples up to a factor $O(n/c^2)$. Any
implementation reads $\Phi$ to compute the mass and draws $\Omega(K)$ samples, so it
does $\Omega(c/\sqrt{np})$ work: at most a quadratic speed-up.
\item \emph{The \#DNNF route.} \Cref{thm:sis} and \cref{prop:dnnf-env} with a DNNF
$E\supseteq C$ (small leaves) read $E$ at every sample and need $\Omega(D(E)/p)$
samples, hence do $\Omega(c/(n\sqrt p))$ work.
\item \emph{Context width.} Let $N$ be an NNF of $C$ as in \cref{thm:cw}, without
frontier gates (threshold gates are expanded by counters, \cref{sec:ctxwidth};
$|N|\le12|C|$ for the gate-by-gate NNF of a $\{\wedge,\vee,\neg,\oplus\}$-circuit). The DNNF that \cref{thm:cw} builds from $N$
has literal and constant leaves, is equivalent to $C$, and has at most
$(|N|+1)(\mathrm{cw}+3)2^{\mathrm{cw}}$ wires, so
$\mathrm{cw}\ge m-14-\log_2\big((n+2)(|N|+1)(\mathrm{cw}+3)/c\big)$.
\item \emph{DFSB.} Every tree built with the exact rules (T), (T$'$), (F), (D), (B) of
\cref{sec:dfsb} in which every leaf $h$ with more than $2n/3$ variables, other than a
refuted one, has $D(h)\ge1/K'$, where $K'<2^{m-6}$, has at least $c\,2^{m-12}/(n+6)^2$ nodes. This
applies to the ideal tree $T^*_{K'}(C)$ when the atom size $a$ is at most $2n/3$, since
its leaves are constants, atoms and (S)-leaves with $D(h)\ge1/K'$. So the completeness
guarantee of \cref{thm:dfsb}, which needs $m_{K'}(C)$ to be at most the node budget,
requires a budget of at least $c^2/(4096(n+6)^2p)$ unless $K'\ge2^{m-6}=c/(64p)$. (Runs
that stop early at (E)-leaves are not covered.)
\end{enumerate}
\end{corollary}

So for random parity checks with $m\le n/6$ and events with $c\ge1/\poly(n)$, the
costs in (i)--(iv) are exponential in $m$: within a polynomial factor of Monte Carlo
for (i), (iii) and (iv), and at least its square root for any use of a DNNF envelope.
For smaller $c$ the bounds weaken; for $c=2^{-m/2}$, for instance, $c^2/p=p^{-1/3}$.
Meanwhile
the presolve computes the factor $2^{-m}$ exactly, and the affine envelope of
$[Hx=s]\wedge R$ with opaque $R$ has $K=1/c$. Conversely, bounded connectivity along a
whole tree suffices.

\begin{proposition}[bounded width; any product distribution]\label{prop:bw-dp}
Let $V=\{Ax=b\}$ and let $\mathcal T$ be a rooted binary tree with leaf set $[n]$ such
that $\lambda_V(X_t)\le w$ for the leaf set $X_t$ of every node $t$. Then $\one_V$ has a
smooth deterministic DNNF, structured by $\mathcal T$, with at most $2n(4^w+1)$ gates;
for every product distribution, $D(V)$ and exact samples from $D\mid V$ are computed
with $O(n4^w)$ arithmetic operations after $\poly(n,m)$ linear algebra; and
$\lambda_{\rm bal}(V)\le w$.
\end{proposition}

The construction (\cref{app:affine}) is a dynamic program over \emph{live partial
syndromes}: at node $t$ the states are the $s=A_{X_t}x_{X_t}$ with $s\in
b+\mathrm{col}(A_{[n]\setminus X_t})$, a coset of $\mathrm{col}(A_{X_t})\cap
\mathrm{col}(A_{[n]\setminus X_t})$, which has $2^{\lambda_V(X_t)}$ elements. Our
implementation uses linear orders (caterpillar trees). A narrow tree can be found by
the fixed-parameter algorithms of Hlin\v{e}n\'y and Oum \cite{hlinenyoum2008} or by
heuristics, or it can be supplied by a prover; the dynamic program is exact for any
tree, and the width only affects time. Let $\mathrm{bw}(V)$ be the least $w$ for
which a tree as in \cref{prop:bw-dp} exists; for $n\ge2$ this is the branch-width of the
binary matroid of $A$ minus one, since the usual connectivity function
\cite{hlinenyoum2008} is $\lambda_V+1$. Together with \cref{thm:parity-barrier}(b) for
$C=E=\one_V$ ($c=1$), the smallest smooth DNNF for $\one_V$ with fan-in-two
$\wedge$-gates and leaves on at most $2n/3$ variables (for instance literal leaves) has
between $2^{\lambda_{\rm bal}(V)}$ and $2n(4^{\mathrm{bw}(V)}+1)$ gates. The two ends are
polynomially related when $\mathrm{bw}(V)=\lambda_{\rm bal}(V)+O(1)$, but not in general: for
$V=V_1\times V_2$ with $V_1,V_2$ on the two halves of $[n]$, $\lambda_{\rm bal}(V)=0$, while
$\mathrm{bw}(V)\ge\mathrm{bw}(V_1)$ can be linear in $n$.

\paragraph{Biased product distributions.}
Linear bijections preserve the uniform measure, but a triangular flip on coordinate $i$
(with $g\not\equiv0$) preserves a product distribution with marginals in $(0,1)$ only if
$q_i=\frac12$. Under a biased product distribution the same tiny circuits become hard.

\begin{theorem}[biased parity events]\label{thm:biased-hard}
Fix a rational $q\in(0,1)\setminus\{\frac12\}$. For $H\in\F_2^{m\times n}$,
$s\in\F_2^m$ and $k\ge1$ let $C$ be the circuit on $N=nk$ inputs $x_{i,c}$ ($i\le n$,
$0\le c<k$)
\[
C=\bigwedge_{i\le n,\;1\le c<k}[x_{i,0}\oplus x_{i,c}=0]\;\wedge\;[Hx_{\cdot,0}=s],
\]
an $\wedge$ of parity checks of total size $O(N+mn)$. Under $D=\mathrm{Bern}(q)^N$,
\[
p=(1-q)^N\sum_{y\in\F_2^n:\,Hy=s}\lambda^{|y|},\qquad\lambda=\Big(\frac q{1-q}\Big)^k.
\]
\begin{enumerate}[label=(\alph*)]
\item If $q<\frac12$ and $k\ge(n+2g+3)/\log_2\frac{1-q}q$, any interval $[L,U]\ni p$
with $U\le2^gL$ decides whether the coset $\{Hy=s\}$ contains a vector of weight at
most $t$, which is $\NP$-complete \cite{bmvt1978}.
\item Consequently, for every fixed $\gamma>0$, a sound procedure that is complete on
these circuits within a factor $2^{N^{1-\gamma}}$ would give $\NP\subseteq\BPP$, and an
argument system (\cref{def:argument}) complete on them within a factor $2$ would give
$\coNP\subseteq\MA$, hence $\PH=\Sigma_2^p$ \cite{bhz1987}.
\end{enumerate}
For $q>\frac12$ complement all inputs, which replaces $s$ by $s+H\one$.
\end{theorem}

The proof is in \cref{app:affine}. The same circuit has $p=2^{-\mathrm{rank}}$ or $0$
under the uniform distribution, and its probability is computed exactly at bounded
width under every product distribution (\cref{prop:bw-dp}). So algebraic tractability
is a property of the circuit and the distribution together. For $s=0$ and all
$q_i\le\frac12$, $p$ is a weight enumerator of $\ker H$, that is, a ferromagnetic
Ising partition function of a binary matroid, which has an FPRAS for graphic
\cite{jerrumsinclair1993} and regular \cite{goldbergjerrum2013regular} matroids and,
under a complexity-theoretic assumption, none for binary matroids in general
\cite{goldbergjerrum2013binary}. In E6 of \cref{tab:affine}, banded LDPC checks
under $\mathrm{Bern}(1/3)$ ($p\approx2.6\cdot10^{-31}$) are estimated to $4\%$ in
$0.2$ seconds by the affine envelope with exact dynamic-programming masses, whereas the
biased-pivot presolve only certifies $p\le1.3\cdot10^{-22}$ within its budget.

\paragraph{Linear thresholds.}
A single linear threshold function (LTF) is tractable under every product
distribution, which makes unions of LTF pieces tractable, but conjunctions are
generally not.

\begin{proposition}[single LTFs \cite{gopalan2011fptas}]\label{prop:ltf}
For $L=[w\cdot x\le\theta]$ with integer weights and a product distribution with
marginals of $\beta$ bits, $D(L)$ has a deterministic FPTAS using $O(n^3\beta/\eps)$ exact
operations on rationals of $\poly(n,\beta,1/\eps)$ bits, and $D\mid L$ has an exact sampler using
$O(n\log(n\beta/\eps))$ expected operations per sample after the same preprocessing.
\end{proposition}

Gopalan, Klivans, Meka, \v{S}tefankovi\v{c}, Vempala and Vigoda note that their
branching-program method counts knapsack solutions under all product and symmetric
distributions \cite[Sec.~3.2]{gopalan2011fptas}. \Cref{app:affine} gives a
self-contained dynamic program, with values rounded to powers of $1+\eps/(2n)$, and a
rejection step that makes the sampler exact. With the rounded masses as certified
lower bounds, the Karp--Luby--Madras estimator remains sound: if
$C=\bigvee_{j\le T}P_j$ where each $P_j$ is an LTF $\wedge$ a subcube, or a
variable-disjoint $\wedge$ of such pieces (and of affine subspaces on uniform
coordinates), the estimator has mean in $[e^{-\eps/3}p,p]$ and looseness at most
$e^{\eps/3}T$ (\cref{app:affine}).

\begin{proposition}[an affine system and one LTF \cite{gopalan2011fptas}]\label{prop:affine-ltf}
Let $C=[Ax=b]\wedge[w\cdot x\le\theta]$ with integer $w$, and let $x_1,\dots,x_n$ be
ordered so that $\lambda_t=\mathrm{rank}A_{\le t}+\mathrm{rank}A_{>t}-\mathrm{rank}A\le
w_0$ for every $t$. Then $D(C)$ has a deterministic FPTAS using
$O(n^32^{w_0}\beta/\eps)$ operations under every product distribution with $\beta$-bit
marginals, and $D\mid C$ has an exact sampler.
\end{proposition}

This is an instance of \cite[Thm.~3.6 with Claim~3.4]{gopalan2011fptas}: for
$V=\{Ax=b\}$, $D\mid V$ is a small-space source of width $2^{w_0}$ in their sense, whose states are the live
syndromes of \cref{prop:bw-dp} and whose transition probabilities are proportional to
$D_t$ times the masses of the suffixes, $[w\cdot x\le\theta]$ is a monotone read-once
branching program in every variable order, and $D(C)=D(V)\Pr_{D\mid V}[w\cdot
x\le\theta]$. Their construction takes time quadratic in the source width. The direct
dynamic program of \cref{app:affine}, which runs the threshold recursion inside each
live syndrome, is linear in $2^{w_0}$ and also gives the exact sampler. Without such
structure, conjunctions are hard.

\begin{proposition}[non-closure]\label{prop:nonclosure}
Under the uniform distribution, deciding $p>0$ is $\NP$-complete for each of the
following classes of polynomial-size circuits: (a) $[w\cdot x\ge\theta]\wedge[w\cdot
x\le\theta]$ (two LTFs; subset sum); (b) $[Hx=s]\wedge[|x|\le t]$ (an affine system and
one LTF; coset weight \cite{bmvt1978}); (c) $[\sum_{uv\in E}(x_u\oplus x_v)\ge k]$ (a
threshold of parities, i.e.\ a quadratic threshold; maximum cut). For (d)
$[|Hx|\le t]$, approximating $p$ within any factor below $2$ is $\NP$-hard
\cite{vardy1997}. Hence no sound procedure is complete on any of these classes unless
$\NP\subseteq\BPP$, and no argument system is complete for upper bounds on any of them
unless $\coNP\subseteq\MA$.
\end{proposition}

\begin{proof}
(a)--(c) are the stated $\NP$-complete problems, and a sound procedure complete on the
class decides them as in \cref{prop:worstcase}(a); an argument system complete for upper
bounds gives $\MA$ proofs of the complements as in \cref{thm:ma}. For (d), $x\mapsto
Hx$ maps the uniform distribution onto the uniform distribution on the code
$\mathcal C=\mathrm{col}(H)$, so $p=|\{y\in\mathcal C:|y|\le t\}|/|\mathcal C|$, which
is $1/|\mathcal C|$ if the minimum distance of $\mathcal C$ exceeds $t$ and at least
$2/|\mathcal C|$ otherwise; deciding which is $\NP$-hard for codes given by generator
matrices \cite{vardy1997} (take $H$ to be the transposed generator matrix).
\end{proof}

The tractable compositions are thus: $\vee$ (Karp--Luby--Madras), variable-disjoint
$\wedge$, $\wedge$ with an affine system whose pivots can be chosen uniform and outside
the other conjunct's support (\cref{prop:presolve}, $p=2^{-r}D(L)$), and affine
$\wedge$ LTF at bounded width (\cref{prop:affine-ltf}).

\paragraph{Fourier and degree structure.}
Rare events with simple Fourier structure are affine.

\begin{proposition}\label{prop:fourier}
Let $D$ be uniform and $C\not\equiv0$.
(a) If $C$ has $k$ nonzero Fourier coefficients then $p\ge1/k$.
(b) If $C$ has degree $d$ as a multilinear polynomial over $\F_2$ or over $\R$ then
$p\ge2^{-d}$.
(c) For every affine $V$ and quadratic polynomial $Q$ over $\F_2$,
$\Pr_{x\in V}[Q(x)=0]\in\{0\}\cup[\frac14,1]$.
\end{proposition}

\begin{proof}
(a) For $x_0$ with $C(x_0)=1$, $1\le\sum_S|\hat C(S)|\le\sqrt{k\sum_S\hat C(S)^2}=
\sqrt{kp}$ by Parseval. (b) We show $\Pr_x[f(x)\ne0]\ge2^{-\deg f}$ for every nonzero multilinear
polynomial $f$ over $\F_2$ or $\R$, by induction on the number of variables, and apply
it to $f=C$. Write $f=x_nG+H$ with $G,H$ free of $x_n$. If $G\equiv0$, then $f=H$ has
fewer variables. Otherwise $\deg G\le\deg f-1$, and wherever $G(x_{<n})\ne0$ at least
one of $f(x_{<n},0)=H$ and $f(x_{<n},1)=H+G$ is nonzero, so $\Pr[f\ne0]\ge\frac12\Pr[G\ne0]\ge
2^{-\deg f}$ (the base case is a nonzero constant). (c) Parametrise $V$ affinely;
$Q$ becomes a quadratic $q$ on $\F_2^{\dim V}$ with associated alternating form
$B(z,z')=q(z+z')+q(z)+q(z')+q(0)$ of rank $2h$. If $h=0$, $q$ is affine and
$\Pr[q=0]\in\{0,\frac12,1\}$. Otherwise
$(\E(-1)^{q})^2=\E_{z'}(-1)^{q(z')+q(0)}\E_z(-1)^{B(z,z')}\le2^{-2h}$, since the inner
expectation vanishes off the radical of $B$; so $\Pr[q=0]=\frac12(1+\E(-1)^q)\in[\frac14,\frac34]$.
\end{proof}

So an event with $p=2^{-50}$ needs Fourier sparsity at least $2^{50}$ and degree at
least $50$, yet $50$ independent parity checks have spectral norm $1$, and Boolean
functions of bounded spectral norm are $\pm$-combinations of boundedly many coset
indicators \cite{greensanders2008}. By (c), one quadratic equation adds at most two
bits of surprise to an affine envelope; affine systems with one quadric are counted
exactly by Dickson's normal form \cite{ek1990}, while systems of quadrics contain SAT.
Sampling on the Fourier side does not help either: for $C=[Ax=b]$ the unbiased
estimator $Y=(-1)^{b\cdot y}\prod_{i:(A^\top y)_i=1}(1-2q_i)$, $y$ uniform in $\F_2^m$,
has $\E Y^2\ge2^{-m}$ (from $y=0$), a relative variance of about $1/p$ when
$p\approx2^{-m}$.

\paragraph{Experiments.}
\Cref{tab:affine} reports our implementation (module \texttt{affine.py}, script
\texttt{exp\_affine.py}; $\eps=0.05$, $\delta=0.01$; wall-clock seconds on a shared
4-core machine). The exact values are
computed without the module: by closed forms for E1--E4 (for E1--E3 the checks
restricted to the inputs outside the residual have full rank, so $p=2^{-m}\Pr[R]$), by
Gaussian elimination on the raw check rows and inclusion--exclusion for E5, by an exact
rational sweep over the bands for E6a, by enumerating the $2^{24}$-element coset for
E6b, and by enumeration in rational arithmetic for E7; the module's dynamic program
agrees with the E6 values to $2\cdot10^{-14}$. Monte Carlo is charged the number of samples its KL interval needs
for the same ratio. E3 illustrates DFSB on a dense parity event at a size where it can
run; with $m=12>n/6$ it lies outside \cref{cor:random-parity}, so it is not an instance
of the barrier. On $36$ inputs with $12$ dense checks ($p=8.8\cdot10^{-5}$), DFSB with a
$400$-node budget spent $79$ seconds and $1.5\cdot10^6$ samples and returned an interval
of ratio $29$, while the presolve needed $1.6\cdot10^4$ samples and $0.02$ seconds; at
this size plain Monte Carlo is cheap too ($3.7\cdot10^7$ bit-parallel evaluations, about
$0.2$ seconds). In $600$ estimates on random $16$-input circuits under uniform and
product distributions (E7), no interval missed the exact value. The whole experiment
takes about $111$ seconds.

The implementation is sound in floating point as well. Threshold gates are evaluated and
propagated exactly, in integer arithmetic after scaling the weights (\cref{sec:setting};
the weights of a repeated child are added), so the threshold rule (X1) of
\cref{sec:xpt} is a valid inference about $C$. Masses are carried as a
mantissa and a binary exponent with proven relative-error bounds (at most
$\gamma_{5n+5}$ for the dynamic program, where $\gamma_k=ku/(1-ku)$ and $u=2^{-53}$),
every interval is rounded outward, and an upper end below the range of doubles becomes
the smallest positive double rather than $0$. The Hoeffding--KL intervals exclude a value
only when its divergence certainly exceeds the bound, allowing for the rounding of the
empirical mean.

\begin{table}[t]
\centering\footnotesize
\setlength{\tabcolsep}{4pt}
\begin{tabular}{@{}lrllrrr@{}}
\toprule
instance (Monte Carlo samples) & $n$ & $p$ & method & $U/L$ & samples & sec.\\
\midrule
E1 gadget checks ($1.2\cdot10^{18}$) & 120 & $1.54\cdot10^{-14}$ & presolve & $1.045$ & $3.3\cdot10^4$ & $0.27$\\
 & & & affine envelope, $K=1.39$ & $1.042$ & $1.6\cdot10^4$ & $0.46$\\
E2 random checks ($4.1\cdot10^{34}$) & 600 & $4.30\cdot10^{-31}$ & presolve & $1.045$ & $3.3\cdot10^4$ & $2.3$\\
 & & & affine envelope, $K=1.38$ & $1.041$ & $1.6\cdot10^4$ & $3.7$\\
E4 Feistel ($4.8\cdot10^{16}$) & 92 & $3.69\cdot10^{-13}$ & triangular presolve & $1.042$ & $6.6\cdot10^4$ & $0.16$\\
 & & & affine envelope & $U{=}5.0\cdot10^{-5}$ & $10^5$ & $0.11$\\
E5 two of three ($3.6\cdot10^{12}$) & 90 & $4.89\cdot10^{-9}$ & affine envelope & $U{=}5.0\cdot10^{-5}$ & $10^5$ & $0.19$\\
 & & & XPT, 7-node argument & $1.040$ & $6.6\cdot10^4$ & $0.43$\\
E6a banded ($6.9\cdot10^{34}$) & 300 & $2.59\cdot10^{-31}$ & affine envelope + DP & $1.041$ & $4.1\cdot10^3$ & $0.24$\\
 & & & biased-pivot presolve & $U{=}1.3\cdot10^{-22}$ & $2\cdot10^5$ & $2.7$\\
E6b copies & 216 & $1.47\cdot10^{-43}$ & affine env., width $\le2^6$ & $U{=}5.1\cdot10^{-5}$ & $2\cdot10^5$ & $0.67$\\
 & & & affine env., width $\le2^{16}$ & $1.006$ & $1.0\cdot10^3$ & $1.9$\\
E3 dense checks ($3.7\cdot10^7$) & 36 & $8.77\cdot10^{-5}$ & presolve & $1.083$ & $1.6\cdot10^4$ & $0.02$\\
 & & & DFSB, 400 nodes & $28.7$ & $1.5\cdot10^6$ & $79$\\
\bottomrule
\end{tabular}
\caption{Affine methods ($\eps=0.05$, $\delta=0.01$; E3: $\eps=0.1$, $\delta=0.05$). E1: 45
checks built from $\wedge,\vee,\neg$ gadgets $\wedge$ a 3-DNF; E2: 100 dense random checks
$\wedge$ a 3-DNF; E4: $[L\oplus f(R)=t]$ with $k=40$ and a nonlinear $f$, $\wedge$ a side
condition; E5: two of three blocks of 14 checks through a threshold gate, $\wedge$ a 3-DNF;
E6: $\mathrm{Bern}(1/3)$ inputs, 100 banded checks (a) and the copies gadget of
\cref{thm:biased-hard} with $n=36$, $m=12$, $k=6$ (b). Every interval contained the exact
$p$; ``$U{=}\dots$'' marks a run that stopped at its sample cap with the valid interval
$[0,U]$.}
\label{tab:affine}
\end{table}

\paragraph{Open questions.} (1) Lower bounds for XOR-DNNF envelopes: an explicit
rare event with $p=2^{-m}$ for which every XOR-DNNF envelope has size times
looseness $2^{\Omega(m)}$. A Fourier version of the coset lemma reduces this to
quadratic systems whose nonzero combinations have large cross-rank across every
balanced linear cut. (2) Is the quadratic gap in \cref{cor:barrier-methods}(i)--(ii)
real? (3) For $s=0$ and $q_i\le\frac12$, $\coNP$-hardness does not apply: can an
$\MA$ prover help, for instance with a certificate of rapid mixing? (4) Detection:
certify affineness of sub-circuits whose intermediate ANFs are dense (parities
computed through adders or lookup tables), for instance from $\F_2$-polynomial-calculus
proofs supplied by a prover.
